% Part: incompleteness
% Chapter: representability-in-q
% Section: comp-representable

\documentclass[../../../include/open-logic-section]{subfiles}

\begin{document}

\olfileid{inc}{req}{crq}
\olsection{توابع محاسبه‌پذیر در~$\Th{Q}$ بازنمایی‌پذیرند}

\begin{thm}
هر تابع محاسبه‌پذیر در~$\Th{Q}$ بازنمایی‌پذیر است.
\end{thm}

\begin{proof}
برای معین‌بودن مطلب و با استفاده از تز چرچ--تورینگ، بگوییم تابعی
محاسبه‌پذیر است اگر و تنها اگر بازگشتی عام باشد. توابع بازگشتی عام
آن‌هایی‌اند که بتوان آن‌ها را از تابع صفر~$\Zero$، تابع
جانشین~$\Succ$ و تابع افکنش~$\Proj{n}{i}$ با استفاده از ترکیب، بازگشت
اولیه و کمینه‌سازی منظم تعریف کرد. بنا بر \olref[pri]{lem:prim-rec}،
هر تابع~$h$ که بتوان آن را از~$f$ و~$g$ تعریف کرد، با استفاده از ترکیب
و کمینه‌سازی منظم نیز از~$f$، $g$ و~$\Zero$، $\Succ$، $\Proj{n}{i}$،
$\Add$، $\Mult$ و~$\Char{=}$ تعریف‌شدنی است. بنابراین تابعی بازگشتی
عام است اگر و تنها اگر بتوان آن را از~$\Zero$، $\Succ$،
$\Proj{n}{i}$، $\Add$، $\Mult$ و~$\Char{=}$ با استفاده از ترکیب و
کمینه‌سازی منظم تعریف کرد.

افزون بر این نشان داده‌ایم که توابع پایهٔ مورد بحث در~$\Th{Q}$
بازنمایی‌پذیرند
(\cref{inc:req:bre:prop:rep-zero,inc:req:bre:prop:rep-succ,inc:req:bre:prop:rep-proj,inc:req:bre:prop:rep-id,inc:req:bre:prop:rep-add,inc:req:bre:prop:rep-mult})،
و هر تابعی که از توابع بازنمایی‌پذیر به‌وسیلهٔ ترکیب یا کمینه‌سازی
منظم تعریف شود
(\olref[cmp]{prop:rep-composition}،
\olref[min]{prop:rep-minimization}) نیز بازنمایی‌پذیر است. پس هر تابع
بازگشتی عام در~$\Th{Q}$ بازنمایی‌پذیر است.
\end{proof}

\begin{explain}
نشان داده‌ایم که مجموعهٔ توابع محاسبه‌پذیر را می‌توان به‌منزلهٔ مجموعهٔ
توابع بازنمایی‌پذیر در~$\Th{Q}$ مشخص کرد. در واقع، اثبات کلی‌تر است.
از تعریف بازنمایی‌پذیری به‌آسانی می‌توان دید هر نظریه‌ای که~$\Th{Q}$
را گسترش دهد (یا بتوان~$\Th{Q}$ را در آن تعبیر کرد) قادر به بازنمایی
توابع محاسبه‌پذیر است. اما برعکس، در هر دستگاه !!{derivation}ی که مفهوم
!!{derivation} در آن محاسبه‌پذیر باشد، هر تابع بازنمایی‌پذیر محاسبه‌پذیر
است. بنابراین، برای نمونه، مجموعهٔ توابع محاسبه‌پذیر را می‌توان
به‌منزلهٔ مجموعهٔ توابع بازنمایی‌پذیر در حساب پئانو، یا حتی نظریهٔ
مجموعه‌های تسرملو--فرنکل، مشخص کرد. چنان‌که گودل اشاره کرد، این امر تا
اندازه‌ای شگفت‌آور است. خواهیم دید که در بحث اثبات‌پذیری، پرسش‌ها به
نظریه‌ای که در نظر می‌گیرید بسیار حساس‌اند؛ به‌طور تقریبی، هرچه اصول
موضوع نیرومندتر باشند، امور بیشتری را می‌توانید ثابت کنید. اما در
طیف گسترده‌ای از نظریه‌های اصل‌موضوعی، توابع بازنمایی‌پذیر دقیقاً همان
توابع محاسبه‌پذیرند؛ نظریه‌های نیرومندتر، مادامی که اصل‌موضوع‌پذیر
باشند، توابع بیشتری را بازنمایی نمی‌کنند.
\end{explain}

\end{document}
